\documentclass[11pt]{article}
\usepackage[UTF8]{ctex}
\usepackage{amsmath,amssymb,amsthm}
\usepackage{etoolbox}

% 1. 自定义定理样式实现标题后强制换行
\newtheoremstyle{break}% 样式名称
  {\topsep}{\topsep}% 环境上下间距
  {\normalfont}% 正文字体
  {}% 缩进量
  {\bfseries}{.}% 头部字体和标点
  {\newline}% 头部后的间距：设置为 \newline 即可实现完美的换行效果
  {}% 头部时间/附加参数规范

% 应用新样式并声明常用的定理类环境
\theoremstyle{break}
\newtheorem{theorem}{定理}[section]
\newtheorem{lemma}[theorem]{引理}
\newtheorem{corollary}[theorem]{推论}
\newtheorem{definition}[theorem]{定义}

% 2. 利用 etoolbox 宏包动态修补 proof 环境，使其标题后自动换行
\patchcmd{\proof}{\item[\hskip\labelsep\itshape#1\@addpunct{.}]}{\item[\hskip\labelsep\itshape#1\@addpunct{.}]\hfill\par}{}{}

\begin{document}
\textbf{固定重数和的集合大小的多项式压缩}

\textbf{草稿}

\textbf{2026年4月29日}

\textbf{摘要}

对于固定的 $h$，令 $N(h,k)$ 为最小的 $N$，使得每一个由 $k$ 元整数集所达到的 $|hA|$ 的基数，也都能被 $\{1,\dots,N\}$ 的某个 $k$ 元子集所达到。我们利用 Rajagopal 关于 $|hA|$ 的可能取值范围的固定-h定理证明：对每个固定的 $h$ 有
\[ N(h,k) \le k^{Ch} \]
成立。新的关键在于用多项式直径的块替换 Rajagopal 的大值构造中的稀疏几何块。该替换在低次希尔伯特函数的层面进行，当 $h$ 固定时，低次希尔伯特函数是唯一与 $h$ 重和集相关的数据。
对于一个有限集 $A \subset \mathbb{Z}$ 和一个正整数 $h$，记
\[ hA = \{a_1 + \cdots + a_h : a_i \in A\} \]
（允许重复）。定义
\[ R(h,k) = \{|hA| : A \subset \mathbb{Z}, |A| = k\}. \]
令 $N(h,k)$ 为满足以下条件的最小正整数 $N$：
\[ R(h,k) = \{|hA| : A \subseteq \{1,\dots,N\}, |A| = k\}. \]
等价地，我们可以在区间 $[0,N-1]$ 上工作，并在最后进行平移。
本文的目的是证明以下多项式压缩界。
\textbf{定理 1.1.} 对每个固定的 $h\geq 1$，存在常数 $C$，使得对所有 $k\geq 2$，有
$$N(h,k) \leq k^{C h}.$$
证明使用了 Rajagopal 关于可能的固定重数和的集合大小的范围的定理。设
$$M_{h,k}=hk-h+1,\quad K_{h,k}=\binom{h+k-1}{h},$$
并定义
$$\Delta_{h,k}=\bigcup_{\ell=0}^{\min(h,k)-3}[M_{h,k}+\ell h+1,\,M_{h,k}+\ell h+(h-2-\ell)].$$
这里及以下，$[u,v]$ 表示整数 $n$ 满足 $u\leq n\leq v$ 的集合；若 $u>v$，则该区间为空。
\end{document}